\centerline{\bf \Huge Домашнее задание.}
\centerline{\bf \Huge Power point}

\begin{flushright}
        \scriptsize{Ты сильнейший потому что ты Годжо Сатору,\\или ты Годжо Сатору потому что ты сильнейший?\\Сугуру Гето.}
\end{flushright}

\problem{1}
В треугольнике $ABC$ высоты $AA_1$, $BB_1$ и $CC_1$ пересекаются в точке $H$. Докажите без использования подобия, что $AH\cdot A_1H = BH \cdot B_1H= CH  \cdot C_1H$. Сформулируйте (без доказательства) обратное утверждение.

\problem{2}
Пусть $I$ --- центр вписанной окружности треугольника $ABC$, $I_a$ --- центр вневписанной окружности, вписанной в угол $A$. Докажите, что $AI \cdot AI_a = AB \cdot AC.$

\problem{3}
Окружность делит каждую из сторон треугольника на $3$ равные части. Докажите с помощью степени точки, что треугольник равносторонний.

\problem{4}
Две окружности таковы,что у них существуют все четыре общие касательные. Докажите, что середины отрезков этих касательных лежат на одной прямой.


\problem{5}
Обозначим за $I$ центр вписанной в треугольники $ABC$ окружности. Перпендикуляр к прямой $AI$, восставленный в точке $I$, пересекает прямую $BC$ в точке $P$ . Точка $Q$ --- основание проекции точки $I$ на прямую $AP$ . Докажите, что четырехугольник $BQAC$ --- вписанный.

\problem{6}
В остроугольном треугольнике из двух вершин провели отрезки к противоположным сторонам. На этих отрезках как на диаметрах построены окружности. Докажите, что общая хорда этих окружностей проходит через ортоцентр треугольника.

\problem{7}
Продолжения сторон $AB$ и $CD$ выпуклого четырехугольника $ABCD$ пересекаются в точке $P$, продолжения сторон $AD$ и $BC$ --- в точке $Q$, а диагонали --- в точке $R$. Рассмотрим окружности, построенные на отрезках $AC$ и $BD$ как на диаметрах. 
\begin{itemize}
    \item[a)] Докажите, что ортоцентр треугольника $APD$ лежит на их радикальной оси.
    \item[b)] \textbf{Прямая Обера}. Докажите, что ортоцентры треугольников $APD, BPC, AQB$ и $DQC$ лежат на одной прямой.
    \item[c)] \textbf{Прямая Гаусса}. Докажите, что середины отрезков $AC, BD$ и $PQ$ лежат на одной прямой.

\end{itemize}